\documentclass[11pt]{article}
\usepackage{amsmath,amssymb,amsthm}
\usepackage[margin=1.1in]{geometry}
\newtheorem{lemma}{Lemma}
\theoremstyle{definition}
\newtheorem{conjecture}{Conjecture}
\theoremstyle{remark}
\newtheorem{remark}{Remark}
\newcommand{\Nbar}{\overline{N}}

\title{The prime content of the Riemann zeros' displacement field:\\
a preregistered measurement programme}
\author{Carl Gribble\thanks{Draft v0.1, companion to Programme Note 4 \cite{Note4} and the
preregistration registry of the code repository \cite{Code}, whose commit hashes date every gate
and prediction below. Developed in collaboration with Claude (Anthropic). Responsibility for the
content rests with the author. A full disclosure of AI use appears at the end of the paper.
\emph{Status:} this draft awaits the full-text literature verification flagged in Section~7 and
the author's review; it should not be circulated as making priority claims before both complete.}}
\date{August 2026 --- draft}

\begin{document}
\maketitle

\begin{abstract}
That the primes are visible in the Riemann zeros is classical: Landau's theorem --- in Gonek's
uniform form, $\sum_{0<\gamma\le T} x^{i\gamma} = -\tfrac{T}{2\pi}\Lambda(x)/\sqrt{x} + E(x,T)$ ---
puts a von Mangoldt spike at every prime power, and Odlyzko's computations made the picture vivid.
This paper measures, under a preregistration discipline borrowed from the experimental sciences,
a quantitative structure inside that classical picture: the per-frequency \emph{amplitude law} with
which the primes enter the zeros' \emph{displacement field} $u_k = \Nbar(\gamma_k) - (k - \tfrac12)$,
the object consumed by canonical-system (Hilbert--P\'olya-flavored) reconstructions. Across four
disjoint windows of the first 500 zeta zeros and the first 150 zeros of $L(s,\chi_{-4})$, with
every gate and prediction committed to a public registry before each first run, we find:
(i) a transfer law $T(\omega) = e^{-a\tau^2}$, $\tau = \omega/2\pi\bar\rho$, with a single constant
$a = 1.15 \pm 0.16$ (block jackknife) that survived two out-of-sample window tests and a
cross-family test at deviations of $0.003$--$0.05$, beating a linear form-factor alternative
$6.6\times$ out of sample; (ii) a localization: the counting function carries the primes at full
Landau--Gonek amplitude at every height tested ($0.99 \pm 0.02$, both families) while only the
zero-sampled displacement sequence attenuates --- so the law is the price of reading the primes from
the zeros' \emph{positions}; and (iii) a marginal-significance overtone structure at prime-power
frequencies (pooled deviations $-0.21 \pm 0.12$ at $m = 9$, $-0.34 \pm 0.22$ at $m = 25$,
$-0.18 \pm 0.30$ at $m = 27$; jointly all negative, individually weak), consistent in kind with the
Bogomolny--Keating lower-order arithmetic terms, including two frequencies committed before their
first measurement anywhere. The family-independence of $a$ is presented as a measured face of
Rudnick--Sarnak universality, not a surprise. One conjecture is stated, with its falsification
protocol and its expected rigorous home (discrete twisted moments of $S$ at the zeros; pair
correlation); the registry records twelve refuted predictions and four instrument errors caught by
frozen gates alongside the confirmations. No progress on the Riemann Hypothesis is claimed.
\end{abstract}

\section{Introduction: what is known, and what is measured here}

Three classical facts frame everything below. First, Landau \cite{Landau} proved, and Gonek
\cite{Gonek} made uniform, that the zeros ring at the primes:
\[
\sum_{0<\gamma\le T} x^{i\gamma} \;=\; -\frac{T}{2\pi}\,\frac{\Lambda(x)}{\sqrt{x}} \;+\; E(x,T),
\]
a von Mangoldt-weighted spike at every prime power and nothing elsewhere. Second, Odlyzko's
large-scale computations \cite{Odlyzko87,OdlyzkoEssay} made the Fourier picture vivid and
established the random-matrix (GUE) statistics of the zeros at short range \cite{Montgomery};
Rudnick and Sarnak \cite{RudnickSarnak} proved the universality of these correlations across
automorphic $L$-functions in restricted regimes; and Bogomolny and Keating \cite{BK1,BK2}
identified the lower-order, explicitly arithmetic corrections that structure Odlyzko's data beneath
the universal terms. Third, on the constructive side, semiclassical potentials whose spectra track
the zeros on average are classical \cite{WuSprung,SchHutch,Berry88,BerryKeating}.

None of the above is at stake in this paper, and none of it is claimed. What this paper adds is a
\emph{measurement}: the quantitative amplitude law with which the primes enter a specific object
--- the zeros' displacement field, $u_k = \Nbar(\gamma_k) - (k - \tfrac12)$ with $\Nbar$ the
Riemann--von Mangoldt smooth counting function --- together with the discipline under which the
measurement was made. The displacement field matters because it is exactly the data consumed by
canonical-system reconstructions of the sought self-adjoint realization (the Hilbert--P\'olya
object): a companion note \cite{Note4} shows that every leading structure of such reconstructions
is arithmetic-free, so whatever the primes contribute to the ``Hamiltonian'' they contribute
through this field.

The findings, in one paragraph. Regressing $u_k$ on $\sin(\gamma_k \ln m)$ over a window with mean
zero density $\bar\rho$ returns the Landau--Gonek amplitude \emph{attenuated}, and the attenuation
follows a Debye--Waller law in the scaling variable $\tau = \omega/2\pi\bar\rho$:
$T(\omega) = e^{-a\tau^2}$ with a single constant across four disjoint zeta windows and a second
$L$-function family. The law is a property of position-sampling alone: measured in continuous $t$
(where $S(t) = N(t) - \Nbar(t)$ is exactly computable between located zeros), the prime content is
full, at every height, in both families --- a numerical verification of Landau--Gonek --- while the
zero-sampled sequence attenuates. On top of the law sits a small, individually marginal but
jointly consistent suppression at prime-power frequencies --- the natural place for the
Bogomolny--Keating arithmetic terms to cast a shadow --- including at two frequencies
($\ln 25$, $\ln 27$) whose measurements were committed before they had ever been made.

\section{Method: the registry}

Every scan below ran under preregistration: instrument constants, gates (PASS/FAIL thresholds),
and predictions were committed to the public repository \cite{Code} \emph{before} each first run;
results are append-only; refuted predictions and instrument errors are permanent. The registry for
this paper's chain comprises seven prereg/results pairs (RH4, RH4b, RH5, RH5b, RH6, RH7, and the
librarian pass), recording along the way: twelve refuted predictions, four instrument or
prereg-level errors caught by frozen gates and fixed with thresholds untouched, one dissolved
analysis plan, and two withheld claims that were only admitted after surviving out-of-sample tests
on data they had never seen. We regard this discipline --- standard in the experimental sciences,
apparently novel in this literature --- as a contribution of the paper independent of any result.

\section{Objects, instruments, and calibration}

Zeros $\gamma_1..\gamma_{500}$ of $\zeta$ (mpmath, 30 digits) and $\gamma_1..\gamma_{150}$ of
$L(s,\chi_{-4})$ (located by a Hardy-function scan with self-verified indexing). Field
$u_k = \Nbar(\gamma_k) - (k-\tfrac12)$; by Riemann--von Mangoldt, $u_k = -\bar S(\gamma_k)$
exactly. Regression of $u_k$ on $\{\sin(\gamma_k \ln m)\}$ over a frequency battery of prime
powers plus polynomial nuisance; measured slopes normalized by the explicit-formula amplitudes
$\Lambda(m)/(\pi\sqrt m \ln m)$ to give ratios $T(m)$; window density $\bar\rho_w$ = mean of
$\ln(\gamma/2\pi)/2\pi$. Windows: $W_1 = $ zeros $1..60$, $W_2 = 141..200$, $W_3 = 301..400$,
$W_4 = 401..500$.

\emph{Calibration against a theorem.} The Landau--Gonek sums over all 500 zeros match their
predicted values at the percent level: measured/predicted $= 1.003, 0.995, 0.966, 0.987, 0.977$
at $r = 2, 3, 4, 5, 9$. The instrument reproduces the classical first-order structure before any
new measurement is read.

\section{The transfer law}

\begin{center}
\begin{tabular}{lcccc}
\hline
window & $m=2$ & $m=3$ & $m=5$ & $m=7$ \\
\hline
$W_1$ (heights 14--163) & $0.921 \pm .026$ & $0.808 \pm .017$ & $0.636 \pm .031$ & $0.550 \pm .061$ \\
$W_2$ (305--396) & $0.984 \pm .020$ & $0.936 \pm .036$ & $0.827 \pm .078$ & $0.735 \pm .115$ \\
$W_3$ (544--680) & $0.974 \pm .016$ & $0.936 \pm .039$ & $0.844 \pm .034$ & $0.825 \pm .052$ \\
$W_4$ (682--811) & $0.973 \pm .032$ & $0.915 \pm .017$ & $0.889 \pm .084$ & $0.772 \pm .059$ \\
\hline
\end{tabular}
\end{center}

\noindent (Ratios with block-jackknife standard errors; the attenuation weakens with height at
every gated frequency, as committed in advance at $W_2$ vs $W_1$.) Two laws in
$\tau = \omega/2\pi\bar\rho_w$ were fitted on $W_1, W_2$ and \emph{recorded before} the $W_3$
test: Debye--Waller $T = e^{-a\tau^2}$ and linear $T = 1 - c\tau$ (the fitted $c = 0.5255$ sits
provocatively near the form-factor value $\tfrac12$). At $W_3$ the Debye--Waller law held with
gated deviations $0.003/0.004/0.036/0.004$ and out-of-sample error $6.6\times$ smaller than the
linear law's; $W_4$ (run later, inside the overtone instrument) added a fourth confirmation at
mean absolute deviation $0.024$. Jackknife over delete-block refits gives
\[
a \;=\; 1.15 \pm 0.16 .
\]
On $L(s, \chi_{-4})$, with the \emph{same} recorded constant: magnitude deviations
$0.022/0.026/0.050$ at $m = 3/5/7$; the even frequencies, forced null by $\chi(2^k) = 0$, measured
silent ($0.001$ at $m = 2$); and the sign pattern is exactly the character's. Given
Rudnick--Sarnak universality \cite{RudnickSarnak}, the family-independence of $a$ is what one
should \emph{expect} if the law derives from pair correlation; we present the cross-family test as
a quantitative face of their theorem, and as evidence for that derivation route, not as a
surprise. A local-variance mechanism ($a$ set by the window's own $\mathrm{Var}\,\bar S$) was
committed twice and refuted twice; the fixed constant transferred better.

\section{Localization: the law is the price of positions}

A structural lemma, machine-checked on synthetic displaced lattices (exact recovery, $1.000$
across nine lines): least squares of a field on its own components \emph{evaluated at the true
sample points} returns full coefficients identically --- sampling displacement cannot attenuate an
in-span line. The attenuation must therefore live elsewhere, and a committed continuous-time
measurement located it: with $S(t) = N(t) - \Nbar(t)$ computed exactly between the located zeros,
the same regression in continuous $t$ returns the \emph{full} Landau--Gonek amplitudes ---
$0.99 \pm 0.02$ across all fifteen cells, three zeta windows and the Dirichlet family, fitted
attenuation constant $\approx 0$ --- while the zero-sampled sequence attenuates as in the table
above. (Our preregistered expectation was the opposite --- that continuous and sampled content
would agree --- and its refutation, recorded in the registry, is precisely what localizes the
law.) The Debye--Waller attenuation is thus a property of reading the primes from the zeros'
\emph{positions}: the counting function carries them in full at every height, as the classical
theory says it must; the displacement sequence --- the data a Hamiltonian reconstruction actually
consumes --- pays the $e^{-a\tau^2}$ price.

\begin{conjecture}[transfer law, localized]
Over a height window $W$ with mean zero density $\bar\rho_W$, the per-frequency content of the
zero-position displacement sequence satisfies
$T(\omega; W) = \exp(-a\,(\omega/2\pi\bar\rho_W)^2)$ with $a = 1.15 \pm 0.16$ independent of
height and family, while the counting fluctuation carries the primes at full explicit-formula
amplitude. \emph{Expected rigorous home:} the object is a discrete twisted moment of $\bar S$ at
the zeros, $\sum_{\gamma} \bar S(\gamma)\, x^{i\gamma}$, within the ambit of the Landau--Gonek
generalization literature and controlled by pair correlation; its GUE universality would explain
the family-independence of $a$. \emph{Falsification protocol:} further windows and families under
the frozen battery; the prime-power frequencies (Section~6); and derivation of the jump-point
sampling weight, which would promote the law toward a theorem.
\end{conjecture}

\section{Overtones at prime powers: a marginal but committed structure}

The frequencies at which measurements departed from the law are exactly the prime powers ---
harmonics of strong prime lines, since $\ln p^k = k \ln p$. Pooled over the four windows, with
jackknife errors: deviations $-0.213 \pm 0.120$ at $m = 9$, $-0.339 \pm 0.221$ at $m = 25$,
$-0.176 \pm 0.304$ at $m = 27$, $-0.036 \pm 0.132$ at $m = 8$, $-0.051 \pm 0.036$ at $m = 4$;
genuine-prime controls $+0.035$ ($m = 11$) and $-0.010$ ($m = 13$). We state the evidential
situation plainly: \emph{individually} these are $0.6$--$1.8\sigma$ effects; the evidence is the
joint sign pattern (all five prime-power deviations negative while both prime controls sit on the
law), the cross-family recurrence of the $m = 9$ suppression, and the fact that $m = 25$ and
$m = 27$ were committed, with thresholds, before either frequency had been measured anywhere ---
both landed negative as predicted. A committed constant-persistence model for $m = 9$ was refuted:
the suppression decays toward the law with height ($0.37/0.39/0.46/0.65$ across the windows), the
natural scaling of a second-order mixing term. This is the place where the Bogomolny--Keating
lower-order arithmetic terms \cite{BK1,BK2} should cast their shadow on the displacement sequence,
and we frame the table as a candidate measured face of their theory, awaiting both more data and a
derivation. The phase-locking input --- Landau sums at $r = 4, 9$, the harmonics' first moments ---
is verified at the percent level (Section~3).

\section{Related work, non-claims, and required verification}

A librarian pass against the live literature (report in \cite{Code}) governs this section. Known
and owned elsewhere: the peaks (Landau; Odlyzko); the continuous full-amplitude content
(Landau--Gonek --- our Section~3 is its numerical verification); correlation universality across
families (Rudnick--Sarnak); arithmetic lower-order structure (Bogomolny--Keating); semiclassical
potentials and their residuals (Wu--Sprung; Berry--Keating; Schumayer--Hutchinson). Claimed here:
the measured amplitude law with its uncertainty, its out-of-sample and cross-family record, the
localization result, the committed overtone table, and the registry discipline. Not claimed:
novelty of any mechanism; any progress on RH. \emph{Required before submission:} full-text
verification against the Landau--Gonek generalization literature and the discrete-moment work of
Fujii \cite{Fujii1,Fujii2} and successors, and against the modern rigorous treatments of Berry's
number-variance theory, none of which could be read in full in the environment where this draft
was prepared; the conjecture's ``we are not aware of a published quantitative law'' stands only
until that reading completes. The de Bruijn--Newman bound of Rodgers--Tao \cite{RodgersTao}
calibrates the programme's ambitions: RH, if true, is barely true, and nothing in this paper
attempts it.

\section{Reproduction}

Every number above is printed by one of seven scripts under \texttt{run\_all.py} in \cite{Code}
(\texttt{coefficient\_primes.py}, \texttt{transfer\_height.py}, \texttt{transfer\_law.py},
\texttt{dirichlet\_transfer.py}, \texttt{transfer\_theory.py}, \texttt{overtone.py},
\texttt{transfer\_errors.py}); the registry directory pairs each instrument with its
preregistration (constants, gates, predictions, all dated by commit hash) and its append-only
results, including every refutation and instrument error mentioned in the text.

\section*{Disclosure of AI use}

This paper and the instruments it reports were developed in an extensive collaboration with
Claude (Anthropic; accessed 2026). The direction of the work and all decisions about scope,
claims, and the preregistration protocol are the author's. Within that direction, the AI system
designed the instruments and preregistrations, committed predictions before runs (several of
which were refuted and are so recorded), performed the analyses and the literature search
reported in Section~7, and drafted the manuscript. Every numerical claim is reproduced by the
cited scripts; the registry records the full prediction-and-refutation history. The author has
reviewed the manuscript and takes full responsibility for its content.

\begin{thebibliography}{99}
\bibitem{Landau} E. Landau, \"Uber die Nullstellen der Zetafunktion, \emph{Math. Ann.} \textbf{71}
(1912), 548--564.
\bibitem{Gonek} S.~M. Gonek, An explicit formula of Landau and its applications to the theory of
the zeta-function, \emph{Contemp. Math.} \textbf{143} (1993), 395--413.
\bibitem{Odlyzko87} A.~M. Odlyzko, On the distribution of spacings between zeros of the zeta
function, \emph{Math. Comp.} \textbf{48} (1987), 273--308.
\bibitem{OdlyzkoEssay} A.~M. Odlyzko, Primes, quantum chaos, and computers, in \emph{Number
Theory}, National Research Council, 1990, 35--46.
\bibitem{Montgomery} H.~L. Montgomery, The pair correlation of zeros of the zeta function, in
\emph{Analytic Number Theory}, Proc. Sympos. Pure Math. XXIV, AMS, 1973, 181--193.
\bibitem{RudnickSarnak} Z. Rudnick and P. Sarnak, Zeros of principal $L$-functions and random
matrix theory, \emph{Duke Math. J.} \textbf{81} (1996), 269--322.
\bibitem{BK1} E.~B. Bogomolny and J.~P. Keating, Gutzwiller's trace formula and spectral
statistics: beyond the diagonal approximation, \emph{Phys. Rev. Lett.} \textbf{77} (1996),
1472--1475.
\bibitem{BK2} E.~B. Bogomolny and J.~P. Keating, Random matrix theory and the Riemann zeros II:
$n$-point correlations, \emph{Nonlinearity} \textbf{9} (1996), 911--935.
\bibitem{Berry88} M.~V. Berry, Semiclassical formula for the number variance of the Riemann
zeros, \emph{Nonlinearity} \textbf{1} (1988), 399--407.
\bibitem{BerryKeating} M.~V. Berry and J.~P. Keating, The Riemann zeros and eigenvalue
asymptotics, \emph{SIAM Rev.} \textbf{41} (1999), 236--266.
\bibitem{WuSprung} H. Wu and D.~W.~L. Sprung, Riemann zeros and a fractal potential, \emph{Phys.
Rev. E} \textbf{48} (1993), 2595--2598.
\bibitem{SchHutch} D. Schumayer and D.~A.~W. Hutchinson, Physics of the Riemann hypothesis,
\emph{Rev. Mod. Phys.} \textbf{83} (2011), 307--330.
\bibitem{Fujii1} A. Fujii, On a theorem of Landau, \emph{Proc. Japan Acad. Ser. A} \textbf{65}
(1989), 51--54.
\bibitem{Fujii2} A. Fujii, On a theorem of Landau, II, \emph{Proc. Japan Acad. Ser. A}
\textbf{66} (1990), 291--296.
\bibitem{RodgersTao} B. Rodgers and T. Tao, The de Bruijn--Newman constant is non-negative,
\emph{Forum Math. Pi} \textbf{8} (2020), e6.
\bibitem{Note4} C. Gribble, The positivity lane: instruments, calibrations, and a requirements
specification for the missing self-adjoint realization, Programme Note 4 (2026), in \cite{Code}.
\bibitem{Code} C. Gribble, Reference implementations and preregistration registry, code
repository (2026), \texttt{https://github.com/carlgribble-caa/prime-moments}.
\end{thebibliography}

\end{document}
